\documentclass[conference]{IEEEtran}
\IEEEoverridecommandlockouts
\usepackage{cite}
\usepackage{amsmath,amssymb,amsfonts}
\usepackage{algorithmic}
\usepackage{graphicx}
\usepackage{textcomp}
\usepackage{xcolor}
\usepackage{listings}
\usepackage[T1]{fontenc}
\usepackage[utf8]{inputenc}
\usepackage{comment}
\usepackage{url}
\usepackage{hyperref}
\usepackage{booktabs}
\usepackage{longtable}
\usepackage{tabularx}
\usepackage{float}

\definecolor{codebg}{rgb}{0.95,0.95,0.95}
\lstset{
  backgroundcolor=\color{codebg},
  basicstyle=\scriptsize\ttfamily,
  keywordstyle=\color{blue},
  commentstyle=\color{gray},
  stringstyle=\color{purple},
  showstringspaces=false,
  breaklines=true,
  frame=single,
  numbers=left,
  numberstyle=\tiny,
  captionpos=b,
  tabsize=2,
  xleftmargin=1em,
}

\def\BibTeX{{\rm B\kern-.05em{\sc i\kern-.025em b}\kern-.08em
    T\kern-.1667em\lower.7ex\hbox{E}\kern-.125emX}}

\begin{document}

\title{LIDAL: Desenvolvimento e Avaliação do Laboratório Interactivo Digital de Álgebra Linear para a Integração Tecnológica no Ensino Superior}

% Bloco de autores construído manualmente com minipages de largura fixa
% e igual (0.48\textwidth cada), em vez do mecanismo automático
% \IEEEauthorblockN/\and do IEEEtran — com 8 autores e afiliações longas,
% esse mecanismo dimensiona cada coluna pelo seu próprio conteúdo e produz
% colunas visivelmente descentradas entre si; com largura fixa e
% \centering em cada minipage, as duas colunas ficam sempre simétricas.
\author{
\begin{minipage}[t]{0.48\textwidth}\centering
1\textsuperscript{st} Manuel Francisco Jamba\\
\textit{Departamento de Ciências Exactas e Naturais}\\
\textit{Instituto Superior de Ciências da Educação da Huíla}\\
Lubango, Angola\\
\href{mailto:m251295@isced-huila.edu.ao}{m251295@isced-huila.edu.ao}
\end{minipage}%
\hfill
\begin{minipage}[t]{0.48\textwidth}\centering
2\textsuperscript{nd} Pedro Mussungo NGalangue\\
\textit{Departamento de Ciências Exactas e Naturais}\\
\textit{Instituto Superior de Ciências da Educação da Huíla}\\
Lubango, Angola\\
\href{mailto:m251300@isced-huila.edu.ao}{m251300@isced-huila.edu.ao}
\end{minipage}\\[1.4em]
\begin{minipage}[t]{0.48\textwidth}\centering
3\textsuperscript{rd} Patrocínio Ndala Kambinda Ntyamba\\
\textit{Departamento de Ciências Exactas e Naturais}\\
\textit{Instituto Superior de Ciências da Educação da Huíla}\\
Lubango, Angola\\
\href{mailto:m251324@isced-huila.edu.ao}{m251324@isced-huila.edu.ao}
\end{minipage}%
\hfill
\begin{minipage}[t]{0.48\textwidth}\centering
4\textsuperscript{th} Mpeheyo Branco Pacheco Himi\\
\textit{Departamento de Ciências Exactas e Naturais}\\
\textit{Instituto Superior de Ciências da Educação da Huíla}\\
Lubango, Angola\\
\href{mailto:m251298@isced-huila.edu.ao}{m251298@isced-huila.edu.ao}
\end{minipage}\\[1.4em]
\begin{minipage}[t]{0.48\textwidth}\centering
5\textsuperscript{th} Nelson de Jesus Henriques Ngonga\\
\textit{Departamento de Ciências Exactas e Naturais}\\
\textit{Instituto Superior de Ciências da Educação da Huíla}\\
Lubango, Angola\\
\href{mailto:m251299@isced-huila.edu.ao}{m251299@isced-huila.edu.ao}
\end{minipage}%
\hfill
\begin{minipage}[t]{0.48\textwidth}\centering
6\textsuperscript{th} Mateus Sacuyema Morais\\
\textit{Departamento de Ciências Exactas e Naturais}\\
\textit{Instituto Superior de Ciências da Educação da Huíla}\\
Lubango, Angola\\
\href{mailto:m251297@isced-huila.edu.ao}{m251297@isced-huila.edu.ao}
\end{minipage}\\[1.4em]
\begin{minipage}[t]{0.48\textwidth}\centering
7\textsuperscript{th} Sávio Panzo Costa Ivo\\
\textit{Departamento de Ciências Exactas e Naturais}\\
\textit{Instituto Superior de Ciências da Educação da Huíla}\\
Lubango, Angola\\
\href{mailto:m251217@isced-huila.edu.ao}{m251217@isced-huila.edu.ao}
\end{minipage}%
\hfill
\begin{minipage}[t]{0.48\textwidth}\centering
8\textsuperscript{th} Marcelino Galvino Thikanha Júlio\\
\textit{Departamento de Ciências Exactas e Naturais}\\
\textit{Instituto Superior de Ciências da Educação da Huíla}\\
Lubango, Angola\\
\href{mailto:m251296@isced-huila.edu.ao}{m251296@isced-huila.edu.ao}
\end{minipage}
}

\maketitle

\begin{abstract}
O ensino da Álgebra Linear no Ensino Superior enfrenta desafios decorrentes do elevado nível de abstracção teórica e da dificuldade dos estudantes em coordenar registos de representação algébricos, numéricos e gráficos. Para responder a esta problemática educacional no contexto do ISCED-Huíla, este artigo apresenta o desenvolvimento e a avaliação empírica do LIDAL (Laboratório Interactivo Digital de Álgebra Linear), uma plataforma web interactiva concebida em Python com os frameworks Streamlit, NumPy, SymPy e Plotly. O ecossistema está organizado em seis módulos matemáticos nucleares, complementados por módulos de exploração gráfica livre, gamificação e conteúdo teórico, e fundamenta-se numa sequência didáctica interactiva baseada no ciclo Explorar--Observar--Interpretar--Experimentar--Verificar. A metodologia seguiu um estudo de caso misto com 32 estudantes do 1.\textordmasculine{} Ano da licenciatura em Ensino da Matemática do ISCED-Huíla. Os resultados quantitativos demonstraram um ganho de aprendizagem estatisticamente significativo do pré-teste para o pós-teste (aumento da média de 42,5\% para 81,2\%, $p < 0{,}001$, $d = 3{,}92$) e uma apreciação altamente favorável nos critérios pedagógicos e de usabilidade (superior a 90\% de concordância). Conclui-se que a mediação tecnológica promovida pelo LIDAL fortalece a intuição geométrica e a capacidade conceptual na resolução de problemas matriciais e vectoriais.
\end{abstract}

\begin{IEEEkeywords}
Álgebra Linear, Laboratório Digital, Ensino Superior, Plotly, Registos de Representação Semiótica, ISCED-Huíla
\end{IEEEkeywords}

\section{Introdução}

\subsection{Contextualização e Problemática}

A integração articulada das Tecnologias de Informação e Comunicação (TIC) no ensino da Matemática constitui uma prioridade pedagógica premente no Ensino Superior contemporâneo. No âmbito do curso de Licenciatura em Ensino da Matemática do Instituto Superior de Ciências da Educação da Huíla (ISCED-Huíla), a unidade curricular de Álgebra Linear desempenha um papel estruturante na formação científica, formal e metodológica dos futuros professores de Matemática.

Contudo, constata-se empiricamente no quotidiano académico que os estudantes do 1.\textordmasculine{} Ano do ISCED-Huíla demonstram acentuadas lacunas na compreensão conceptual dos conteúdos programáticos. As dificuldades não residem prioritariamente na memorização mecânica de regras operatórias, mas sim no estabelecimento de pontes reflexivas entre as definições algébricas formais, as manipulações numéricas procedimentais e as suas correspondentes interpretações geométricas no plano e no espaço.

As dificuldades manifestam-se de forma crítica em tópicos fundamentais como: a) Operações algébricas com matrizes de ordem elevada e interpretação do produto matricial; b) Determinação de determinantes e compreensão da sua relação com o factor de escala de áreas e volumes; c) Condições de existência e cálculo da matriz inversa $A^{-1}$; d) Discussão, classificação e representação geométrica da intersecção de subespaços em sistemas de equações lineares ($AX = B$); e) Operações vectoriais em $\mathbb{R}^2$ e $\mathbb{R}^3$, tais como o produto escalar, produto vectorial e combinações lineares; f) Determinação da equação característica $\det(A-\lambda I) = 0$, autovalores e autovectores. Frequentemente, os estudantes encaram o determinante de uma matriz ou a matriz inversa como meros algoritmos aritméticos desprovidos de significado visual, desconhecendo a sua aplicação em transformações geométricas de bases e espaços vectoriais.

\subsection{Questão de Investigação e Hipóteses}

Perante este diagnóstico educacional no contexto do ISCED-Huíla, formula-se a seguinte pergunta de investigação: \textit{De que forma um laboratório digital interactivo, desenvolvido em Python, pode contribuir para a compreensão conceptual e para a motivação na aprendizagem dos conteúdos de Álgebra Linear, através da articulação dinâmica entre representações algébricas, numéricas e gráficas nos estudantes do 1.\textordmasculine{} Ano do ISCED-Huíla?}

Como hipóteses de investigação, definem-se:

\textbf{$H_0$ (Hipótese Nula)}: A utilização do laboratório digital LIDAL não produz alterações estatisticamente significativas no desempenho académico dos estudantes entre o pré-teste e o pós-teste.

\textbf{$H_1$ (Hipótese Alternativa)}: A mediação tecnológica promovida pelo LIDAL proporciona um aumento estatisticamente significativo no desempenho conceptual e operacional dos estudantes, acompanhado de elevada satisfação técnica e pedagógica.

\subsection{Objectivos do Trabalho}

Para responder à pergunta de investigação, o presente trabalho fixou como objectivo geral conceber, implementar e avaliar o LIDAL (Laboratório Interactivo Digital de Álgebra Linear) como recurso tecnológico e pedagógico de apoio ao ensino e à aprendizagem de Álgebra Linear no ISCED-Huíla.

Os objectivos específicos alinham-se da seguinte forma: 1) Implementar módulos computacionais interaccionais para operações fundamentais com matrizes e vectores em $\mathbb{R}^2$ e $\mathbb{R}^3$; 2) Desenvolver ferramentas dinâmicas para o cálculo simbólico e numérico de determinantes e matrizes inversas com verificação de etapas intermédias; 3) Proporcionar a visualização e resolução de sistemas de equações lineares através do método matricial $AX = B$ e representação gráfica de rectas e planos; 4) Construir simuladores gráficos para a transformação linear associada a valores e vectores próprios; 5) Validar empiricamente a usabilidade e a eficácia pedagógica do sistema com estudantes do ISCED-Huíla.

O público-alvo primário do LIDAL são estudantes do 1.\textordmasculine{} ano de cursos de formação inicial de professores de Matemática no Ensino Superior, com particular incidência nos estudantes do curso de Licenciatura em Ensino da Matemática do ISCED-Huíla. Secundariamente, o sistema destina-se também a docentes de Álgebra Linear que pretendam utilizá-lo como recurso de apoio à leccionação, quer em sala de aula (\textit{modo turma}), quer como material de estudo autónomo recomendado aos estudantes.

\subsection{Importância e Contribuição}

As principais contribuições deste projecto residem na disponibilização de um protótipo funcional de software educacional, tanto em versão web como em executável autónomo para uso offline, na formulação de uma sequência didáctica guiada em 5 etapas com \textit{feedback} imediato, e na produção de evidências empíricas sobre a integração das tecnologias no Ensino Superior em Angola.

\section{Fundamentação Teórica}

\subsection{Teoria dos Registos de Representação Semiótica}

A arquitectura pedagógica do LIDAL apoia-se na Teoria dos Registos de Representação Semiótica desenvolvida por Raymond Duval \cite{duval1995}. Duval estabelece que os objectos matemáticos não são apreensíveis pelos sentidos directos, existindo exclusivamente através das suas representações semióticas (símbolos, tabelas, gráficos, linguagem natural). Consequentemente, a aprendizagem matemática cognitiva exige o domínio de dois processos fundamentais: 1) \textit{Tratamento}: Transformação de uma representação noutra dentro do mesmo registo semiótico (por exemplo, aplicar a regra de Sarrus para calcular um determinante no registo numérico); 2) \textit{Conversão}: Transformação de uma representação num registo para outro registo semioticamente distinto (por exemplo, converter a equação de um sistema linear na intersecção gráfica de dois planos no espaço $\mathbb{R}^3$).

O LIDAL foi concebido para actuar como um catalisador de \textit{conversões}, permitindo que a alteração de um coeficiente no registo algébrico altere instantaneamente a renderização no registo gráfico.

\subsection{Os Três Mundos da Matemática de David Tall}

Complementarmente, a teoria dos Três Mundos da Matemática proposta por David Tall \cite{tall2004} descreve o desenvolvimento do pensamento matemático através de três domínios relacionais:

\begin{enumerate}
  \item \textbf{Mundo Embodied} (Incorporado / Visuo-espacial): Baseado na percepção sensorial, intuição e visualização espacial dos objectos;
  \item \textbf{Mundo Symbolic} (Simbólico / Procedimental): Baseado na utilização de algoritmos, manipulação de símbolos e cálculos;
  \item \textbf{Mundo Formal} (Axiomático): Baseado em definições rigorosas, teoremas e demonstrações lógicas dedutivas.
\end{enumerate}

O LIDAL estabelece a transição contínua entre o mundo incorporado (gráfico) e o mundo simbólico (cálculo numérico/simbólico), pavimentando o caminho para a abstracção formal.

\subsection{Enquadramento no Modelo TPACK}

O desenvolvimento da ferramenta alinha-se igualmente com o modelo TPACK (\textit{Technological Pedagogical Content Knowledge}), ao integrar o Conhecimento Tecnológico (programação Python, Streamlit), o Conhecimento Pedagógico (sequência didáctica orientada e \textit{feedback}) e o Conhecimento do Conteúdo (conteúdos programáticos de Álgebra Linear).

\subsection{Gamificação e Feedback Imediato na Aprendizagem}

O módulo de Jogos e Desafios do LIDAL fundamenta-se em dois corpos teóricos complementares. Primeiro, o conceito de \textit{gamificação}, definido por Deterding \textit{et al.} \cite{deterding2011} como a utilização de elementos de \textit{design} de jogos (pontuação, progressão por níveis, desafios de dificuldade crescente) em contextos não lúdicos, com o objectivo de aumentar o envolvimento e a motivação intrínseca dos utilizadores --- princípio operacionalizado nos três modos de jogo do sistema (cronometrado, livre por tópico, e desafios com pistas progressivas). Segundo, a revisão seminal de Hattie e Timperley \cite{hattie2007} sobre o impacto do \textit{feedback} na aprendizagem, que identifica o \textit{feedback} imediato e orientado ao processo (não apenas ao resultado final) como um dos factores de maior efeito no desempenho académico. Esta segunda referência fundamenta directamente o desenho do modo ``Desafio com Pistas'', em que o sistema não revela de imediato se a matriz introduzida tem ou não inversa, solicitando antes o cálculo intermédio do determinante e disponibilizando uma pista processual antes de confirmar a resposta final.

\subsection{Análise Comparativa com Softwares Existentes}

A literatura internacional corrobora as dificuldades estruturais no ensino da Álgebra Linear. Moro \textit{et al.} \cite{moro2016} destacam que a abordagem excessivamente formal afasta os alunos da intuição prática. Singh \textit{et al.} \cite{singh2021} analisaram 20 softwares educacionais voltados para a Álgebra Linear, concluindo que a maioria actua como calculadoras numéricas estáticas desprovidas de acompanhamento didáctico.

Softwares consolidados como o GeoGebra \cite{hohenwarter2004} oferecem recursos genéricos potentes, mas exigem elevado esforço de programação de applets por parte do docente. Plataformas recentes como o CG.lab \cite{damiano2025} restringem-se à Álgebra Linear focada na Computação Gráfica, enquanto o repositório LAVIMAT \cite{araujo2024} reúne links dispersos sem coerência algorítmica unificada. O LIDAL diferencia-se por combinar rigor formal, simulação interactiva e desafios gamificados integrados no currículo universitário (módulo de Jogos e Desafios, descrito na Secção~\ref{sec:modulos-complementares}).

\begin{table*}[htbp]
\caption{Análise Comparativa Detalhada entre o LIDAL e Plataformas Educacionais da Literatura}
\label{tab:comparativa}
\centering
\begin{tabular}{|l|l|l|l|l|}
\hline
\textbf{Dimensão Analisada} & \textbf{GeoGebra \cite{hohenwarter2004}} & \textbf{LAVIMAT \cite{araujo2024}} & \textbf{CG.lab \cite{damiano2025}} & \textbf{LIDAL (Proposto)} \\
\hline
Foco do Conteúdo & Geral (K-12 ao Ensino Superior) & Repositório Genérico & Computação Gráfica & Álgebra Linear Universitária \\
\hline
Registos Semióticos & Simbólico e Gráfico & Texto/Hipertexto & Gráfico 2D/3D & Algébrico, Numérico e Gráfico \\
\hline
Sequência Didáctica & Não Estruturada & Passiva & Exercícios Fixos & Guiada em 5 Etapas + Desafios \\
\hline
Feedback ao Aluno & Apresentação Visual & Ausente & Verificação Certo/Errado & Passo a Passo Explicativo \\
\hline
Processamento Simbólico & CAS Parcial & Não Possui & Não Possui & Integrado via SymPy (Python) \\
\hline
Arquitectura Tecnológica & Java / HTML5 Engine & Web / CMS Links & Next.js / React Three & Python / Streamlit / Plotly \\
\hline
\end{tabular}
\end{table*}

\section{Metodologia}

\subsection{Desenho do Estudo}

Adoptou-se um desenho de investigação baseado no Estudo de Caso Misto (qualitativo e quantitativo) com aplicação de Pré-Teste e Pós-Teste num único grupo experimental submetido a intervenção pedagógica com o LIDAL.

\subsection{Caracterização dos Participantes}

A amostra do estudo foi constituída por $N = 32$ estudantes (18 do sexo masculino e 14 do sexo feminino, com idades compreendidas entre os 18 e os 24 anos) do 1.\textordmasculine{} Ano do curso de Licenciatura em Ensino da Matemática do ISCED-Huíla, no decorrer do ano lectivo de 2026. A participação foi voluntária e antecedida da assinatura do Termo de Consentimento Livre e Esclarecido.

\subsection{Procedimento e Sequência Didáctica}
\label{sec:seq-didactica}

A intervenção pedagógica decorreu durante 4 semanas (totalizando 16 horas lectivas de laboratório) e seguiu a Sequência Didáctica em 5 Etapas:

1) \textbf{Explorar}: Entrada livre de matrizes e vectores pelos estudantes para geração de hipóteses; 2) \textbf{Observar}: Análise visual imediata dos gráficos e resultados numéricos no ecrã; 3) \textbf{Interpretar}: Leitura do passo a passo explicativo e do embasamento teórico formal; 4) \textbf{Experimentar}: Teste de situações limite (ex.: matrizes singulares, vectores ortogonais); 5) \textbf{Verificar}: Resolução de desafios propostos pelo sistema com validação em tempo real.

\subsection{Metodologia de Desenvolvimento do Sistema}

Do ponto de vista da engenharia de software, o LIDAL foi construído seguindo uma abordagem incremental: cada funcionalidade foi implementada, integrada e validada isoladamente antes de se avançar para a seguinte, em vez de se deixar a integração para uma fase final única. O código-fonte foi mantido sob controlo de versões (Git), com um histórico granular de alterações, e cada alteração funcional foi acompanhada da execução da suite de testes automatizados descrita na Secção~\ref{sec:testes}, funcionando como mecanismo de validação contínua e prevenindo a introdução de regressões ao longo do período de desenvolvimento de 15 dias.

\section{Concepção e Implementação}

\subsection{Arquitectura do Sistema}

O LIDAL foi desenvolvido adoptando uma arquitectura em camadas separadas, garantindo alta modularidade, facilidade de manutenção e rapidez de execução:

\begin{itemize}
  \item \textbf{Camada de Interface de Utilizador (UI)}: Desenvolvida com a biblioteca Streamlit, proporcionando um ambiente web responsivo, intuitivo e com renderização reactiva em tempo real, com navegação por página dedicada a cada módulo funcional;
  \item \textbf{Camada do Núcleo Computacional Matemático}: Desenvolvida em Python 3.12, combinando a biblioteca NumPy para operações matriciais numéricas vectorizadas e a biblioteca SymPy para cálculo algébrico simbólico, determinantes exactos e simplificação de equações formais;
  \item \textbf{Camada de Renderização Gráfica}: Construída sobre a biblioteca Plotly, responsável pela geração dinâmica e interactiva dos vectores no plano e espaço, planos de sistemas lineares e autovectores, incluindo animação de parâmetros por controlo deslizante (\textit{slider}) e reprodução automática (\textit{Play}), executada inteiramente no navegador do utilizador sem pedidos adicionais ao servidor por cada fotograma.
\end{itemize}

O sistema encontra-se disponível tanto como aplicação web quanto como executável autónomo para Windows (empacotado com PyInstaller, biblioteca \textit{onefile}), permitindo a sua utilização offline no ISCED-Huíla sem necessidade de ligação permanente à internet.

\subsection{Justificação das Tecnologias Utilizadas}

A lista de tecnologias sugeridas para este projecto incluía Python, NumPy, SymPy, Matplotlib ou Jupyter Notebook. O LIDAL adopta as três primeiras, mas substitui o par Matplotlib/Jupyter Notebook por Streamlit e Plotly, por razões concretas relacionadas com o público-alvo e com o objectivo pedagógico do sistema. Um caderno Jupyter Notebook, embora adequado para exploração individual por um utilizador já familiarizado com Python, exige que o estudante instale e opere um ambiente de programação --- uma barreira de entrada desadequada ao público-alvo primário (estudantes do 1.\textordmasculine{} ano, muitos sem experiência prévia de programação). O Streamlit resolve esta limitação ao compilar o mesmo código Python numa interface web imediatamente utilizável, sem qualquer conhecimento de programação por parte do utilizador final, mantendo simultaneamente o código-fonte inteiramente em Python (sem exigir HTML, CSS ou JavaScript por parte da equipa de desenvolvimento).

A escolha do Plotly em detrimento do Matplotlib prende-se especificamente com a natureza interactiva exigida pela sequência didáctica \textit{Explorar--Observar--Interpretar--Experimentar--Verificar} (Secção~\ref{sec:seq-didactica}): o Matplotlib produz figuras estáticas, exigindo um novo pedido ao servidor Streamlit (e um novo cálculo em Python) sempre que o utilizador altera um parâmetro; o Plotly, pelo contrário, permite pré-calcular no servidor um conjunto de fotogramas (\textit{frames}) correspondentes a uma gama de valores de um parâmetro, e entregar ao navegador um controlo deslizante e um botão de reprodução automática que animam esses fotogramas inteiramente no lado do cliente, sem qualquer novo pedido ao servidor. Esta arquitectura foi decisiva para tornar viável, sem latência perceptível, a exploração livre de parâmetros pedida na etapa \textit{Explorar} da sequência didáctica --- razão pela qual se prevê, como trabalho futuro (Secção~\ref{sec:conclusao}), construir também uma versão equivalente em Matplotlib, unicamente para fins de comparação técnica objectiva entre as duas abordagens.

\subsection{Módulos Funcionais e Formalismo Matemático}

\subsubsection{Módulo M1 --- Matrizes}
Permite a definição de matrizes $A \in \mathbb{M}_{m \times n}(\mathbb{R})$ e $B \in \mathbb{M}_{p \times q}(\mathbb{R})$. Implementa a validação da condição de conformidade de dimensões para a soma ($m=p \wedge n=q$) e para a multiplicação ($n=p$), calculando o produto segundo a fórmula formal:
\begin{equation}
c_{ij} = \sum_{k=1}^{n} a_{ik}b_{kj}, \quad i = 1,\dots,m;\; j = 1,\dots,q
\label{eq:produto}
\end{equation}

\subsubsection{Módulo M2 --- Determinantes / Inversa (Cálculo do Determinante)}
Os Módulos M2 e M3 correspondem, na aplicação real, a uma única página --- \textit{Determinantes / Inversa} --- apresentada aqui em duas vertentes matemáticas distintas por clareza expositiva. Calcula o determinante $\det(A)$ para matrizes quadradas de ordem $n \geq 2$:
\begin{equation}
\det(A) = \sum_{j=1}^{n} (-1)^{i+j} a_{ij} \det(M_{ij})
\label{eq:determinante}
\end{equation}
Além do valor numérico exacto obtido via SymPy, o sistema renderiza graficamente o paralelogramo formado pelos vectores linha/coluna em $\mathbb{R}^2$, demonstrando que $|\det(A)|$ corresponde à área da figura transformada.

\subsubsection{Módulo M3 --- Determinantes / Inversa (Matriz Inversa)}
Verifica a condição de não-singularidade ($\det(A) \neq 0$). Caso a matriz seja invertível, calcula a matriz adjunta $\mathrm{adj}(A)$ e apresenta a matriz inversa $A^{-1}$:
\begin{equation}
A^{-1} = \frac{1}{\det(A)} \mathrm{adj}(A)
\label{eq:inversa}
\end{equation}
O módulo inclui a verificação algébrica automática do produto $A \cdot A^{-1} = I_n$.

\subsubsection{Módulo M4 --- Sistemas Lineares}
Suporta a resolução do sistema matricial $AX = B$ através da eliminação aplicada à matriz aumentada $[A \mid B]$:
\begin{equation}
AX = B \iff
\begin{bmatrix} a_{11} & \dots & a_{1n} \\ \vdots & \ddots & \vdots \\ a_{m1} & \dots & a_{mn} \end{bmatrix}
\begin{bmatrix} x_1 \\ \vdots \\ x_n \end{bmatrix} =
\begin{bmatrix} b_1 \\ \vdots \\ b_m \end{bmatrix}
\label{eq:sistema}
\end{equation}
Paralelamente, gera a representação geométrica 2D (rectas concorrentes, paralelas ou coincidentes) ou 3D (intersecção de planos), associando a solução algébrica à posição relativa dos subespaços.

\subsubsection{Módulo M5 --- Vetores}
Permite a inserção de vectores $u, v \in \mathbb{R}^2$ ou $\mathbb{R}^3$. Calcula o módulo $\|u\| = \sqrt{\sum u_i^2}$, o produto escalar:
\begin{equation}
u \cdot v = \|u\|\|v\|\cos(\theta) = \sum_{i=1}^{n} u_i v_i
\label{eq:escalar}
\end{equation}
e gera a representação gráfica vectorial da soma $u+v$ através da regra do paralelogramo, além do produto externo e combinações lineares $c_1 u + c_2 v$.

\subsubsection{Módulo M6 --- Valores/Vetores Próprios}
Determina os autovalores $\lambda \in \mathbb{R}$ resolvendo o polinómio característico:
\begin{equation}
p(\lambda) = \det(A - \lambda I) = 0
\label{eq:caracteristico}
\end{equation}
Para cada autovalor obtido, resolve o sistema homogéneo $(A-\lambda I)v = 0$ para encontrar o espaço próprio associado $E_\lambda$. Graficamente, o módulo exibe como a transformação $T(v) = Av$ actua sobre o espaço, destacando que os autovectores mantêm a sua direcção original sofrendo apenas um estiramento/compressão pelo factor $\lambda$.

\subsubsection{Módulos Complementares}
\label{sec:modulos-complementares}
Além dos seis módulos matemáticos nucleares, o LIDAL integra três módulos complementares que operacionalizam de forma transversal a etapa \textit{Verificar} da sequência didáctica: um módulo de \textbf{Exploração Gráfica} de equações livres (estilo GeoGebra, com detecção automática de parâmetros e animação por controlo deslizante); um módulo de \textbf{Jogos e Desafios}, com três modos de gamificação --- jogo cronometrado, modo livre por tópico e desafios com pistas progressivas (o sistema solicita primeiro o cálculo do determinante, com pista disponível, antes de perguntar se a matriz tem inversa) --- e registo de pontuação partilhado por turma; e um módulo de \textbf{Conteúdo}, com fundamentação teórica e exemplos resolvidos por tópico.

\begin{table}[htbp]
\caption{Resumo dos Módulos Funcionais do LIDAL e Recursos Gráficos}
\label{tab:modulos}
\centering
\footnotesize
\begin{tabular}{|p{0.5cm}|p{1.9cm}|p{3.6cm}|}
\hline
\textbf{Mód.} & \textbf{Domínio Matemático} & \textbf{Recurso de Visualização} \\
\hline
M1 & Operações Matriciais & Representação tabular e validação visual \\
\hline
M2 & Determinantes & Área do paralelogramo em $\mathbb{R}^2$ \\
\hline
M3 & Inversão Matricial & Transformação de base unitária \\
\hline
M4 & Sistemas $AX=B$ & Intersecção gráfica de rectas e planos \\
\hline
M5 & Geometria Vectorial & Vectores no espaço $\mathbb{R}^2/\mathbb{R}^3$ e paralelos \\
\hline
M6 & Autoespaços & Rectas de direcção própria e estiramento \\
\hline
\end{tabular}
\end{table}

\subsection{Implementação de Código Principal}

O Código~\ref{lst:autovalores} ilustra a implementação simplificada do motor computacional do Módulo M6 em Python, utilizando SymPy para cálculo simbólico e Plotly para plotagem gráfica interactiva.

\begin{lstlisting}[language=Python, caption={Trecho de código Python do núcleo do Módulo M6 no LIDAL.}, label={lst:autovalores}]
import streamlit as st
import numpy as np
import sympy as sp
import plotly.graph_objects as go

def calcular_autovalores_gui(matriz_data):
    # Definicao da matriz simbolica via SymPy
    A = sp.Matrix(matriz_data)
    lam = sp.symbols('lambda')

    # Polinomio caracteristico det(A - lambda*I)
    poly = A.charpoly(lam)
    autovectores = A.eigenvects()

    st.write("### Polinomio Caracteristico:")
    st.latex(f"p(\\lambda) = {sp.latex(poly.as_expr())} = 0")

    # Renderizacao grafica interativa com Plotly
    fig = go.Figure()
    A_num = np.array(matriz_data, dtype=float)

    for val, mult, vects in autovectores:
        for v in vects:
            v_num = np.array(v, dtype=float).flatten()
            tv = A_num @ v_num  # T(v) = A.v

            fig.add_trace(go.Scatter(
                x=[0, v_num[0]], y=[0, v_num[1]],
                mode="lines+markers",
                name=f"v (lambda={float(val):.2f})"))
            fig.add_trace(go.Scatter(
                x=[0, tv[0]], y=[0, tv[1]],
                mode="lines+markers",
                line=dict(dash="dash"), name="T(v)"))

    fig.update_xaxes(range=[-5, 5])
    fig.update_yaxes(range=[-5, 5], scaleanchor="x")
    st.plotly_chart(fig, width="stretch")
\end{lstlisting}

\section{Resultados e Avaliação}

\subsection{Testes e Verificação Funcional do Sistema}
\label{sec:testes}

Previamente à avaliação com estudantes, o sistema foi submetido a um processo de verificação funcional automatizada, através de uma suite de aproximadamente 40 scripts de teste construídos com o \textit{framework} \texttt{streamlit.testing.v1.AppTest}. Os testes cobrem a totalidade dos módulos funcionais (M1--M6) e complementares, verificando: a ausência de excepções em tempo de execução em cada página da aplicação; a correcção dos resultados numéricos e simbólicos face a casos de referência conhecidos (por exemplo, determinantes, matrizes inversas e valores próprios de matrizes-padrão com solução analítica trivial); e o comportamento correcto da interface perante entradas inválidas ou limite (dimensões incompatíveis, matrizes singulares, sistemas indeterminados ou impossíveis). Esta suite foi executada após cada alteração ao longo do desenvolvimento incremental do sistema, funcionando como mecanismo de regressão contínua e garantindo a estabilidade funcional do produto entregue.

\subsection{Exemplos de Funcionamento}

A título ilustrativo, apresentam-se dois exemplos reais de utilização do sistema. No Módulo M4, introduzindo o sistema de equações $x+y=3$, $x-y=1$ (exemplo pré-carregado por omissão no módulo de Sistemas Lineares), o LIDAL devolve, em simultâneo, a solução exacta $(x,y)=(2,1)$ obtida por eliminação sobre a matriz aumentada, a lista de passos intermédios do cálculo, e o gráfico interactivo das duas rectas com o ponto de intersecção assinalado, permitindo ao estudante verificar visualmente que a solução algébrica corresponde ao ponto de cruzamento geométrico. No Módulo M6, para a matriz diagonal $A = \begin{bmatrix} 2 & 0 \\ 0 & 3 \end{bmatrix}$ (exemplo pré-carregado por omissão no módulo de Valores/Vetores Próprios), o sistema calcula os valores próprios $\lambda_1=2$ e $\lambda_2=3$ com as respectivas direcções próprias alinhadas com os eixos coordenados, e gera uma animação interactiva que interpola entre a matriz identidade e $A$ (parâmetro $t \in [0,1]$ em $M(t)=(1-t)I+tA$), tornando visível o efeito de escala associado a cada valor próprio à medida que a transformação se constrói.

\subsection{Desempenho Académico (Aprendizagem)}

A Tabela~\ref{tab:desempenho} apresenta a evolução do desempenho dos estudantes entre o Pré-Teste e o Pós-Teste, discriminada por módulo conceptual.

\begin{table}[htbp]
\caption{Comparação do Desempenho Académico por Módulo (Média $\pm$ DP em \%)}
\label{tab:desempenho}
\centering
\footnotesize
\begin{tabular}{|p{2.3cm}|c|c|c|}
\hline
\textbf{Módulo} & \textbf{Pré} & \textbf{Pós} & \textbf{Ganho} \\
\hline
M1 --- Op. Matriciais & $54{,}2 \pm 12{,}1$ & $88{,}5 \pm 6{,}4$ & +63,3\% \\
\hline
M2 --- Determinantes & $41{,}8 \pm 14{,}3$ & $79{,}4 \pm 8{,}1$ & +89,9\% \\
\hline
M3 --- Matriz Inversa & $38{,}5 \pm 15{,}2$ & $82{,}1 \pm 7{,}5$ & +113,2\% \\
\hline
M4 --- Sistemas Lineares & $46{,}0 \pm 13{,}8$ & $84{,}6 \pm 8{,}2$ & +83,9\% \\
\hline
M5 --- Vetores $\mathbb{R}^2/\mathbb{R}^3$ & $48{,}1 \pm 11{,}5$ & $86{,}3 \pm 6{,}9$ & +79,4\% \\
\hline
M6 --- Autovalores/vetores & $26{,}4 \pm 10{,}2$ & $66{,}3 \pm 9{,}8$ & +151,1\% \\
\hline
\textbf{Média Global} & $\mathbf{42{,}5 \pm 11{,}2}$ & $\mathbf{81{,}2 \pm 8{,}4}$ & \textbf{+91,1\%} \\
\hline
\end{tabular}
\end{table}

A aplicação do teste $t$ de Student para amostras emparelhadas revelou uma diferença estatisticamente altamente significativa entre a média do pré-teste e do pós-teste ($t(31) = 15{,}68$, $p < 0{,}001$).
\begin{equation}
t = \frac{\bar{d}}{s_d / \sqrt{N}}
\label{eq:teste-t}
\end{equation}

O tamanho do efeito calculado através do $d$ de Cohen resultou em $d = 3{,}92$, indicando um impacto pedagógico de magnitude extremamente elevada:
\begin{equation}
d = \frac{\bar{X}_1 - \bar{X}_2}{s_{\text{pooled}}}
\label{eq:cohen}
\end{equation}

\subsection{Avaliação de Usabilidade (U1--U6)}

Os resultados do questionário Likert referente à usabilidade técnica (U1 a U6) encontram-se resumidos na Tabela~\ref{tab:usabilidade}.

\begin{table}[htbp]
\caption{Resultados da Avaliação de Usabilidade do Sistema ($N=32$)}
\label{tab:usabilidade}
\centering
\footnotesize
\begin{tabular}{|p{3.6cm}|p{0.85cm}|p{0.75cm}|p{0.85cm}|}
\hline
\textbf{Indicador de Usabilidade} & \textbf{Concorda} & \textbf{Neutro} & \textbf{Discorda} \\
\hline
U1. Interface intuitiva e fácil navegação & 93,8\% & 6,2\% & 0,0\% \\
\hline
U2. Acesso rápido às funcionalidades & 87,5\% & 12,5\% & 0,0\% \\
\hline
U3. Clareza das instruções no ecrã & 90,6\% & 6,3\% & 3,1\% \\
\hline
U4. Estética e atractivo visual & 84,4\% & 15,6\% & 0,0\% \\
\hline
U5. Estabilidade operacional do software & 81,3\% & 12,5\% & 6,2\% \\
\hline
U6. Rapidez no tempo de resposta & 93,8\% & 6,2\% & 0,0\% \\
\hline
\end{tabular}
\end{table}

\subsection{Avaliação dos Critérios Pedagógicos (P1--P6)}

A dimensão pedagógica (P1 a P6) foi avaliada de forma excepcionalmente positiva, com 100\% de concordância na capacidade do software em conectar a álgebra à geometria, como exposto na Tabela~\ref{tab:pedagogico}.

\begin{table}[htbp]
\caption{Resultados da Avaliação dos Critérios Pedagógicos ($N=32$)}
\label{tab:pedagogico}
\centering
\footnotesize
\begin{tabular}{|p{3.6cm}|p{0.85cm}|p{0.75cm}|p{0.85cm}|}
\hline
\textbf{Indicador Pedagógico} & \textbf{Concorda} & \textbf{Neutro} & \textbf{Discorda} \\
\hline
P1. Auxílio na compreensão formal & 96,9\% & 3,1\% & 0,0\% \\
\hline
P2. Aumento da motivação para aprender & 90,6\% & 9,4\% & 0,0\% \\
\hline
P3. Articulação Álgebra/Geometria & 100,0\% & 0,0\% & 0,0\% \\
\hline
P4. Integração Teoria--Prática & 93,8\% & 6,2\% & 0,0\% \\
\hline
P5. Utilidade do feedback em tempo real & 93,8\% & 6,2\% & 0,0\% \\
\hline
P6. Relevância dos recursos visuais & 100,0\% & 0,0\% & 0,0\% \\
\hline
\end{tabular}
\end{table}

\section{Discussão}

Os achados empíricos deste estudo confirmam a viabilidade das teorias de Raymond Duval \cite{duval1995} e David Tall \cite{tall2004} na mediação tecnológica do ensino superior. Ao preencher a lacuna entre os algoritmos numéricos e a intuição espacial, o LIDAL permitiu aos estudantes superar os obstáculos cognitivos associados à abstracção dos espaços vectoriais relatados por Moro \textit{et al.} \cite{moro2016}.

A comparação com as conclusões de Singh \textit{et al.} \cite{singh2021} reforça que softwares educacionais não devem limitar-se a fornecer respostas prontas; a inclusão de uma sequência didáctica guiada com \textit{feedback} em tempo real revelou-se determinante para a consolidação da aprendizagem.

\subsection{Potencialidades}

Para além dos ganhos de aprendizagem quantificados na Secção~V-C, os resultados evidenciam quatro potencialidades concretas do sistema. Primeiro, a articulação instantânea entre os registos algébrico, numérico e gráfico (Tabela~V, indicador P3: 100\% de concordância) confirma empiricamente, no contexto do ISCED-Huíla, o papel do LIDAL como catalisador de \textit{conversões} semióticas na acepção de Duval \cite{duval1995}. Segundo, a verificação algébrica automática integrada nos Módulos M3 ($A \cdot A^{-1} = I_n$) e M6 ($Av = \lambda v$) torna visível ao estudante, sem esforço adicional, a coerência entre o resultado obtido e a definição formal subjacente --- respondendo directamente à dificuldade, diagnosticada na Secção~I-A, de os estudantes tratarem estas operações como algoritmos sem significado. Terceiro, a arquitectura de fotogramas pré-calculados do Plotly (Secção~IV-B) permite explorar em tempo real toda a gama de um parâmetro sem latência perceptível, tornando praticável, mesmo em ligações mais lentas, a etapa \textit{Explorar} da sequência didáctica. Quarto, a disponibilidade simultânea como aplicação web e como executável autónomo alarga o alcance do sistema a contextos com acesso à internet limitado ou instável, um factor pedagogicamente relevante no panorama educativo angolano.

\subsection{Limitações}

Não obstante os resultados favoráveis, a solução apresenta limitações que devem ser reconhecidas. Do ponto de vista metodológico, o desenho de investigação assentou num único grupo experimental, sem grupo de controlo, e numa amostra reduzida ($N=32$) circunscrita a uma única instituição, o que limita a generalização dos resultados obtidos; estudos futuros deverão contemplar amostras mais amplas e desenhos comparativos com grupo de controlo. Do ponto de vista técnico, a visualização gráfica interactiva dos Módulos M2, M3 e M6 encontra-se actualmente limitada a matrizes de ordem 2 e 3, não sendo suportada a representação geométrica de matrizes de ordem superior, ainda que o cálculo simbólico e numérico permaneça válido para qualquer ordem. O registo de pontuação do modo colaborativo (``modo turma'') do módulo de Jogos e Desafios, quando o sistema é utilizado através do executável autónomo, não é persistido entre execuções, estando limitado à sessão de utilização corrente. Adicionalmente, os tipos de letra utilizados na interface dependem de uma ligação à internet activa, degradando-se de forma graciosa, mas sem preservar integralmente a identidade visual, em contexto totalmente offline. Estas limitações não comprometem a validade dos resultados obtidos, mas delimitam o âmbito de generalização das conclusões e constituem oportunidades concretas de desenvolvimento futuro.

\section{Conclusão e Trabalhos Futuros}
\label{sec:conclusao}

O projecto LIDAL cumpriu com distinção todos os objectivos traçados, demonstrando que a integração de ferramentas dinâmicas em Python no Ensino Superior promove ganhos significativos no rendimento dos estudantes.

Como trabalhos futuros, pretende-se expandir o laboratório com suporte a óculos de realidade aumentada em 3D, um módulo de \textit{Learning Analytics} para o docente, uma versão de comparação da camada gráfica em Matplotlib para fins de avaliação técnica, e a disponibilização do código-fonte num repositório público sob licença de código aberto.

\nocite{gladcheff2001,oliveira2001,bardin2016}
\bibliographystyle{IEEEtran}
\bibliography{bibliografia}

\end{document}
